\documentclass[11pt,a4paper]{article}
\usepackage[T1]{fontenc}
\usepackage{lmodern}
\usepackage[margin=27mm,headheight=14pt]{geometry}
\usepackage{amsmath,amssymb,amsthm,mathtools}
\usepackage{microtype,booktabs,longtable,array}
\usepackage{xcolor,xurl}
\definecolor{linkink}{RGB}{28,61,91}
\usepackage[colorlinks=true,allcolors=linkink,pdfusetitle]{hyperref}
\usepackage{fancyhdr}
\pagestyle{fancy}
\fancyhf{}\fancyhead[L]{\small Resonant polylogarithm transseries}\fancyhead[R]{\small ProveIt research draft}
\fancyfoot[C]{\thepage}
\newtheorem{theorem}{Theorem}[section]
\newtheorem{proposition}[theorem]{Proposition}
\newtheorem{corollary}[theorem]{Corollary}
\newtheorem{lemma}[theorem]{Lemma}
\theoremstyle{definition}\newtheorem{definition}[theorem]{Definition}
\theoremstyle{remark}\newtheorem{remark}[theorem]{Remark}
\newcommand{\Li}{\operatorname{Li}}
\newcommand{\PV}{\operatorname{PV}}
\newcommand{\Res}{\operatorname{Res}}
\newcommand{\Borel}{\mathcal B}
\newcommand{\Bell}{\mathcal Y}
\newcommand{\dd}{\,\mathrm d}
\newcommand{\ii}{\mathrm i}
\newcommand{\ee}{\mathrm e}
\newcommand{\N}{\mathbb N}
\newcommand{\C}{\mathbb C}
\newcommand{\Q}{\mathbb Q}
\newcommand{\R}{\mathbb R}
\newcommand{\coeff}[2]{[#1^{#2}]}
\DeclareMathOperator{\sinc}{sinc}
\numberwithin{equation}{section}
\allowdisplaybreaks[2]
\title{Resonant Polylogarithm Transseries:\\Gamma Borel Kernels, Cyclotomic Order Jets,\\and Sharp Dilogarithm Reversion}
\author{Research draft prepared for Vladimir Reshetnikov\\\small OpenAI ChatGPT}
\date{11 October 2026}
\begin{document}
\maketitle
\begin{abstract}
We study two operations that expose information invisible in the terminated
large-argument expansion of an integer-order polylogarithm: differentiation
in its order and nonlinear inversion. Starting from the classical complete
Fermi--Dirac expansion, we construct an explicit lattice Borel kernel and
identify the exponentially small completion with its branch signs fixed.
At every positive integer order, the first order derivative has a meromorphic
Borel kernel expressed by a subtracted Gamma product. Its residues recover an
entire small-argument polylogarithm. A periodic-weight extension gives finite
digamma and cyclotomic-polylogarithm coordinates, without discarding the zero
Fourier mode. All higher order-derivative coefficients are given by harmonic
Bell polynomials, and an exact Gamma quotient resolves the joint large-order
and integer-resonance limit.

The main reversion result concerns $F_2(x)=-\Li_2(-\ee^x)$. Its inverse about
its exact quadratic core has alternating exponential sectors whose
coefficients are positive rational Laurent polynomials. We prove their exact
radius of convergence, identify the unique dominant square-root singularity,
and obtain the explicit $N^{-3/2}$ coefficient law. On the physical curve,
the sector series converges if and only if the target value is at least
$\pi^2/4$; convergence at equality is absolute. The proof extends to a
quadratic polylogarithm-feedback family. The article separates these
derived results from classical input, supplies conventional proofs and
reproducible checks, and records the limits of the source audit.
\end{abstract}
\medskip
\noindent\textbf{Status.} Research draft with explicit proofs; not externally
refereed and not Lean-formalized. The sharp reversion package and the
order-jet/cyclotomic synthesis are proposed contributions, not certified
claims of worldwide priority. Numerical diagnostics are not interval proofs.
\newpage
\tableofcontents
\newpage

\section{Research setting and the main connection}\label{sec:introduction}

The two repository areas meet naturally at a distinction between a formal
asymptotic expansion and the analytic function that realizes it. In the
normalization used here,
\begin{equation}\label{eq:Fdef}
 F_s(x):=-\Li_s(-\ee^x),\qquad
 T_s(x):=F_s(-x)=-\Li_s(-\ee^{-x}),\qquad x>0.
\end{equation}
The subscript is the \emph{polylogarithm order}. The conventional normalized
Fermi--Dirac index is $s-1$, not $s$. In particular, the inverse studied below
is the inverse in $x$ of $-\Li_2(-\ee^x)$, not a choice of inverse branch in
the argument of $\Li_2$ itself.

At positive integer order $m$, the algebraic asymptotic expansion terminates:
\begin{equation}\label{eq:intro-polynomial}
 P_m(x)=\frac{x^m}{m!}
 +\sum_{k=1}^{\lfloor m/2\rfloor}
       \frac{2\eta(2k)}{(m-2k)!}x^{m-2k},
 \qquad \eta(z)=(1-2^{1-z})\zeta(z).
\end{equation}
Nevertheless,
\begin{equation}\label{eq:intro-exact}
 F_m(x)=P_m(x)-(-1)^mT_m(x).
\end{equation}
Thus termination does not erase exponentially small information. Moreover,
although the late coefficients vanish at $s=m$, their derivatives in $s$
do not. Differentiation reveals a factorially divergent tail whose Borel
poles carry the same small-argument polylogarithm. Inversion, meanwhile,
converts the small tail into infinitely many coupled exponential sectors.
The paper makes both mechanisms explicit.

\subsection{What is classical and what is proposed here}

The integral representation and functional equations for polylogarithms
are classical; we use the conventions in DLMF~\cite{DLMFpoly}.
Complete Fermi--Dirac expansions including exponentially small terms predate
this article: see Garoni--Frankel--Glasser~\cite{GFG} and the explicit account
in Gil--Segura--Temme, Section~3 and Lemma~3.1~\cite{GST}. We do \emph{not}
claim the complete Sommerfeld formula or its cosecant coefficient generator
as new.

Likewise, contour Lagrange inversion, Rouch\'e certification, and analytic
square-root transfer are established methods. The repository already has a
careful nonlinear Stokes-transport theorem, including countably summable
perturbations, in \cite{RepoTransport}. We reuse that architecture rather
than renaming it as a new general theorem. The coefficient-transfer step is
classical singularity analysis~\cite{FO}.

The proposed additions are the assembled, model-specific statements:
\emph{the subtracted Gamma kernel for resonant order jets; its finite
periodic-weight digamma dictionary; the all-order harmonic coefficient
formula and detuning limit; and especially the positive Laurent formula,
exact radius, unique dominant singularity, and sharp physical threshold for
the dilogarithm inverse}. The latter conclusions require more than a formal
Lagrange formula or a sufficient small-perturbation disk.

The current polylogarithm corpus already distinguishes accepted identities,
experimental evidence, and pending analytic claims~\cite{RepoPoly}.
The transseries corpus also contains polylogarithmic feedback with a
\emph{linear} dominant core and nonanalytic boundary criticality
\cite{RepoCritical}. Our quadratic feedback theorem is an interior-fold
result with a different core. It is not presented as the first connection
between polylogarithms and transseries, nor as a solution of all of that
corpus's boundary problems.

\subsection{A preview of the sharp threshold}

Let $y>\pi^2/6$ and put
\begin{equation}\label{eq:u-y}
 u=\sqrt{2y-\pi^2/3}.
\end{equation}
We prove that the inverse has the sector representation
\begin{equation}\label{eq:intro-inverse}
 F_2^{-1}(y)=u+\sum_{N\ge1}(-1)^{N-1}a_N(u)\ee^{-Nu},
 \qquad a_N(u)\in\Q_{>0}[u^{-1}],
\end{equation}
precisely when $y\ge\pi^2/4$, including absolute convergence at equality.
For $\pi^2/6<y<\pi^2/4$, the analytic inverse still exists, but the terms of
this particular sector series do not tend to zero.

This threshold is a theorem about convergence of an \emph{exact-core
transseries chart}. It is not a claim about nonexistence of the inverse,
a numerical failure threshold for all algorithms, or a mere optimal
truncation rule. Its proof occupies Sections~\ref{sec:dilog}--\ref{sec:physical}.

\subsection{Conventions and proof dependencies}

The logarithm is real on $(0,\infty)$ and is continued with explicit lateral
signs when necessary. A superscript $+$ on a Borel integral means the path
passes \emph{above} the positive Borel axis. Consequently, for $\xi>n$, the
factor $1-(\xi+\ii0)/n$ has argument $-\pi$. This sign matters.

All asymptotic statements in the sector index $N$ hold with the positive
core parameter $u$ fixed unless uniformity is explicitly stated. All
power-series identities are identities of analytic germs in their indicated
variables. The finite exact computations in the package check specified
instances; the all-order assertions rest on the proofs below.

\section{The classical completion in lattice Borel coordinates}\label{sec:borel}

\subsection{An exact split of the defining integral}

For $\Re s>0$ the standard Fermi--Dirac representation is
\begin{equation}\label{eq:FDintegral}
 F_s(x)=\frac1{\Gamma(s)}
       \int_0^\infty\frac{t^{s-1}}{1+\ee^{t-x}}\dd t.
\end{equation}
It immediately gives the following convenient exact decomposition.

\begin{lemma}[Thermal split]\label{lem:split}
For $x>0$ and $\Re s>0$,
\begin{equation}\label{eq:split}
 F_s(x)=\frac{x^s}{\Gamma(s+1)}+
 \frac1{\Gamma(s)}\int_0^\infty
 \frac{(x+u)^{s-1}-\mathbf1_{\{u<x\}}(x-u)^{s-1}}
      {\ee^u+1}\dd u.
\end{equation}
The identity and its order derivatives are locally uniform in the half-plane
$\Re s>0$.
\end{lemma}
\begin{proof}
Split \eqref{eq:FDintegral} at $t=x$. On $t<x$ use
$(1+\ee^{t-x})^{-1}=1-(1+\ee^{x-t})^{-1}$ and then set $u=x-t$.
On $t>x$ set $u=t-x$. The integral of $t^{s-1}$ from $0$ to $x$ is
$x^s/s$. This gives the displayed formula. On a compact order set with
$\Re s\ge\epsilon>0$, endpoint singularities are bounded by a constant times
$v^{\epsilon-1}(1+|\log v|^r)$ for each fixed derivative order $r$;
the upper endpoint has exponential decay. These bounds justify order
differentiation and local uniformity.
\end{proof}

Expanding the odd difference near $u=0$ yields the formal Sommerfeld series
\begin{equation}\label{eq:Sommerfeld}
 \widehat P_s(x)=\frac{x^s}{\Gamma(s+1)}+
 \sum_{k\ge1}\frac{2\eta(2k)}{\Gamma(s-2k+1)}x^{s-2k}.
\end{equation}
The coefficient identity uses
$\int_0^\infty u^{2k-1}/(\ee^u+1)\dd u=\Gamma(2k)\eta(2k)$.
We do not interchange an infinite, potentially divergent, expansion with
\eqref{eq:split}; finite Taylor expansion with its remainder, or the Borel
construction below, supplies the asymptotic interpretation.

\subsection{A normally convergent Borel lattice}

Remove the leading power by writing
\begin{equation}\label{eq:Hformal}
 \widehat H_s(x)=1+2\sum_{k\ge1}\eta(2k)
                    s^{\underline{2k}}x^{-2k}.
\end{equation}
Our Borel-minor convention is
\begin{equation}\label{eq:Bconvention}
 \Borel\!\left(\sum_{n\ge1}h_nx^{-n}\right)(\xi)
       =\sum_{n\ge1}\frac{h_n}{(n-1)!}\xi^{n-1}.
\end{equation}

\begin{theorem}[Lattice kernel and exact lateral completion]\label{thm:lattice}
The Borel minor of $\widehat H_s-1$ is
\begin{equation}\label{eq:bs}
 b_s(\xi)=s\sum_{n\ge1}\frac{(-1)^{n-1}}{n}
       \left[\left(1+\frac\xi n\right)^{s-1}
             -\left(1-\frac\xi n\right)^{s-1}\right].
\end{equation}
This sum converges normally on compact subsets of any coherent continued
sheet away from the possible singularities $\pm1,\pm2,\ldots$. For
$\Re s>0$, the lateral Laplace sums
\begin{equation}\label{eq:Ps}
 P_s^\pm(x)=\frac{x^s}{\Gamma(s+1)}
       \left(1+\int_{0}^{\infty\ee^{\pm\ii0}}
                         \ee^{-x\xi}b_s(\xi)\dd\xi\right)
\end{equation}
satisfy
\begin{align}
 P_s^\pm(x)&=F_s(x)+\ee^{\mp\ii\pi s}T_s(x),\label{eq:completion}\\
 P_s^+(x)-P_s^-(x)&=-2\ii\sin(\pi s)T_s(x).\label{eq:jump}
\end{align}
Thus, for the arithmetic average $P_s^{\mathrm{av}}=(P_s^++P_s^-)/2$,
\begin{equation}\label{eq:average}
 F_s=P_s^{\mathrm{av}}-\cos(\pi s)T_s.
\end{equation}
All these statements can be differentiated any fixed number of times in
$s$, locally uniformly in $\Re s>0$.
\end{theorem}
\begin{proof}
For $|\xi|<1$, expand the two binomials. The coefficient of $\xi^{2k-1}$
is
\[
 2s\binom{s-1}{2k-1}\sum_{n\ge1}\frac{(-1)^{n-1}}{n^{2k}}
 =\frac{2\eta(2k)s^{\underline{2k}}}{(2k-1)!},
\]
as required. On a fixed compact set, the bracket in the $n$th term is
$O(n^{-1})$ as $n\to\infty$, uniformly in $s$ on compact sets. The extra
factor $1/n$ proves normal convergence. Only finitely many terms require
nontrivial branch choices along any fixed compact continuation path.

On rays strictly above or below the positive axis, the kernel has at most
polynomial--logarithmic growth, uniformly for $s$ in a compact subset of
$\Re s>0$. To see this, split the lattice sum at $n\asymp|\xi|$;
for larger $n$ the preceding difference estimate applies, while the finitely
many smaller terms are bounded by their power growth on the ray. The
same estimates with logarithmic factors hold for order derivatives.
The lateral boundary powers are locally integrable because $\Re s>0$.
They define \eqref{eq:Ps}, equivalently by ray deformation and boundary
limits.

In the $n$th integral make the substitution $\xi=nu/x$. Absolute
integrability of the large-$n$ tail, and local integrability of the finitely
many branch points, justify summing the integrals. The geometric sum is
$\sum_{n\ge1}(-1)^{n-1}\ee^{-nu}=(\ee^u+1)^{-1}$. Consequently the integral
in \eqref{eq:Ps}, after its prefactor is included, becomes
\[
 \frac1{\Gamma(s)}\int_0^\infty
 \frac{(x+u)^{s-1}-(x-u)_{\mp}^{s-1}}{\ee^u+1}\dd u.
\]
For $u>x$, the last power is
$\ee^{\mp\ii\pi(s-1)}(u-x)^{s-1}$. Comparing with \eqref{eq:split} gives
\[
 P_s^\pm-F_s
 =-\frac{\ee^{\mp\ii\pi(s-1)}}{\Gamma(s)}
       \int_x^\infty\frac{(u-x)^{s-1}}{\ee^u+1}\dd u
 =\ee^{\mp\ii\pi s}T_s.
\]
This proves the signs, the jump, and the average. The same compact-order
majorants used in Lemma~\ref{lem:split} justify differentiation.
Finally, the Laplace integral of a finite Taylor remainder of the germ
$b_s$ gives \eqref{eq:Sommerfeld} asymptotically; the part away from the
origin is exponentially small on a fixed admissible ray.
\end{proof}

\begin{corollary}[Integer resonance]\label{cor:integer}
For every positive integer $m$, $b_m$ is a polynomial, both lateral sums
are the polynomial $P_m$ of \eqref{eq:intro-polynomial}, and
\begin{equation}\label{eq:integer-exact}
 F_m(x)=P_m(x)+(-1)^m\Li_m(-\ee^{-x}).
\end{equation}
The perturbative Stokes jump vanishes, but the physical exponential
completion generally does not.
\end{corollary}
\begin{proof}
Each binomial in \eqref{eq:bs} is a polynomial at $s=m$.
Equivalently, $1/\Gamma(m-2k+1)=0$ when $2k>m$.
Equation~\eqref{eq:completion} at $s=m$ now proves the identity.
\end{proof}

For example,
\begin{align*}
 F_1(x)&=x+\log(1+\ee^{-x}),\\
 F_2(x)&=\frac{x^2}{2}+\frac{\pi^2}{6}+\Li_2(-\ee^{-x}),\\
 F_3(x)&=\frac{x^3}{6}+\frac{\pi^2x}{6}-\Li_3(-\ee^{-x}).
\end{align*}
These checks fix the parity convention independently of Borel terminology.

\section{Integer order derivatives and a Gamma Borel kernel}\label{sec:jets}

\subsection{The factorial tail restored by differentiation}

Write
\begin{equation}\label{eq:ak}
 A_k(s)=\frac{2\eta(2k)}{\Gamma(s-2k+1)}.
\end{equation}
If $n=2k-m\ge1$, Gamma reflection gives the exact local identity
\begin{equation}\label{eq:gamma-detuning}
 A_k(m+\epsilon)=2(-1)^{m+1}\eta(2k)
           \frac{\sin(\pi\epsilon)}\pi\Gamma(n-\epsilon).
\end{equation}
In particular,
\begin{equation}\label{eq:first-late}
 A_k(m)=0,\qquad A_k'(m)=2(-1)^{m+1}\eta(2k)\Gamma(2k-m).
\end{equation}
Thus the first order derivative has infinitely many nonzero late
coefficients even though the undifferentiated series terminates.

Define
\begin{align}\label{eq:Km}
 K_m(x)={}&\frac{x^m}{m!}\bigl(\log x-\psi(m+1)\bigr)\notag\\
 &+\sum_{k=1}^{\lfloor m/2\rfloor}
 \frac{2\eta(2k)x^{m-2k}}{(m-2k)!}
       \bigl(\log x-\psi(m-2k+1)\bigr).
\end{align}
This is the finite power--log part of the first order derivative.
Put $\tau_0=1$ and $\tau_{2k}=2\eta(2k)$ for $k\ge1$.

\begin{theorem}[Subtracted Gamma kernel]\label{thm:Gamma-kernel}
For $m\ge1$ the ordinary Borel minor of the late part of
$\partial_s\widehat P_s|_{s=m}$ is
\begin{align}
 B_m(\xi)
 &=(-1)^{m+1}\xi^{-m-1}
   \left\{\Gamma(1+\xi)\Gamma(1-\xi)
           -\sum_{k=0}^{\lfloor m/2\rfloor}\tau_{2k}\xi^{2k}\right\}
       \label{eq:Gamma-kernel}\\
 &=(-1)^{m+1}\xi^{-m-1}
   \left\{\frac{\pi\xi}{\sin\pi\xi}
           -\sum_{k=0}^{\lfloor m/2\rfloor}\tau_{2k}\xi^{2k}\right\}.
       \label{eq:csc-kernel}
\end{align}
The apparent singularity at zero is removable. Its positive poles have
residues
\begin{equation}\label{eq:residues}
 \Res_{\xi=n}B_m(\xi)=\frac{(-1)^{m+n+1}}{n^m},\qquad n\ge1.
\end{equation}
For $J_m(x)=\partial_sF_s(x)|_{s=m}$, one has the exact identity
\begin{equation}\label{eq:jet-PV}
 J_m(x)=K_m(x)+\PV\int_0^\infty\ee^{-x\xi}B_m(\xi)\dd\xi
                   -(-1)^m\left.\partial_sT_s(x)\right|_{s=m}.
\end{equation}
Here the principal value is the arithmetic average of the two lateral
integrals, equivalently the sum of symmetric local principal values at the
positive integer poles. In particular,
\begin{equation}\label{eq:jet-jump}
 \int_+\ee^{-x\xi}B_m(\xi)\dd\xi
 -\int_-\ee^{-x\xi}B_m(\xi)\dd\xi
 =-2\pi\ii(-1)^mT_m(x).
\end{equation}
\end{theorem}
\begin{proof}
The tail in \eqref{eq:first-late} has Borel minor
\begin{equation}\label{eq:Bm-series}
 2(-1)^{m+1}\sum_{k>m/2}\eta(2k)\xi^{2k-m-1}.
\end{equation}
The classical generating identity
$\pi\xi/\sin\pi\xi=1+2\sum_{k\ge1}\eta(2k)\xi^{2k}$
and Gamma reflection~\cite{DLMFgamma} give
\eqref{eq:Gamma-kernel}--\eqref{eq:csc-kernel}. The first exponent in
\eqref{eq:Bm-series} is nonnegative. The residue follows from
$(\sin\pi\xi)'|_{\xi=n}=\pi(-1)^n$.

For the exact, rather than only formal, assertion, let
$c_m=\partial_s b_s|_{s=m}$. Differentiating the lattice kernel $m$ times
in $\xi$ gives
\begin{equation}\label{eq:Bm-partial}
 \frac{c_m^{(m)}(\xi)}{m!}
 =\sum_{n\ge1}\frac{(-1)^{n-1}}{n^m}
      \left(\frac1{n+\xi}-\frac{(-1)^m}{n-\xi}\right)=B_m(\xi).
\end{equation}
Indeed, the $m$th derivative of
$(1+a\xi)^{m-1}\log(1+a\xi)$ is
$(m-1)!a^m/(1+a\xi)$; polynomial terms disappear.
The series converges normally away from its poles.

On either tilted Laplace ray, repeated integration by parts yields
\[
 x^m\int\ee^{-x\xi}c_m(\xi)\dd\xi
 =\int\ee^{-x\xi}c_m^{(m)}(\xi)\dd\xi
       +\sum_{j=0}^{m-1}x^{m-1-j}c_m^{(j)}(0).
\]
The polynomial endpoint contributions, combined with the derivative of
$x^s/\Gamma(s+1)$ in \eqref{eq:Ps}, are exactly $K_m$.
For verification, the endpoint at $j=2k-1<m$ contributes
\[
 \frac{2\eta(2k)x^{m-2k}}{(m-2k)!}
       \bigl(\psi(m+1)-\psi(m-2k+1)\bigr),
\]
which changes the common $\log x-\psi(m+1)$ to the term in
\eqref{eq:Km}. Therefore
\begin{equation}\label{eq:jet-lateral}
 \left.\partial_sP_s^\pm(x)\right|_{s=m}
     =K_m(x)+\int_\pm\ee^{-x\xi}B_m(\xi)\dd\xi.
\end{equation}
Taking the order derivative of \eqref{eq:completion}, and then averaging,
proves \eqref{eq:jet-PV}. Subtraction proves \eqref{eq:jet-jump}.
The sums of pole contributions are absolutely convergent after multiplication
by $\ee^{-nx}$, so the indicated global principal value is well defined.
\end{proof}

The first two kernels illustrate the different local parity:
\begin{align}
 B_1(\xi)&=\frac{\pi}{\xi\sin\pi\xi}-\frac1{\xi^2},\label{eq:B1}\\
 B_2(\xi)&=-\frac1{\xi^3}
       \left(\frac{\pi\xi}{\sin\pi\xi}-1-\frac{\pi^2\xi^2}{6}\right).
\end{align}
For the upper lateral integral,
\begin{equation}\label{eq:upper-jet-real}
 J_m=K_m+\int_+\ee^{-x\xi}B_m(\xi)\dd\xi
           -(-1)^m\bigl(T_m'-\ii\pi T_m\bigr),
\end{equation}
where $T_m'$ denotes an \emph{order derivative}, not an $x$ derivative.
Equation~\eqref{eq:upper-jet-real} is the form used in the independent
complex-ray numerical checks.

\subsection{All higher order coefficients in harmonic coordinates}

The complete exponential Bell polynomials are specified by
\begin{equation}\label{eq:Bell-def}
 \exp\left(\sum_{j\ge1}b_j\frac{z^j}{j!}\right)
       =\sum_{r\ge0}\Bell_r(b_1,\ldots,b_r)\frac{z^r}{r!},
 \qquad \Bell_0=1.
\end{equation}
Let $H_n^{(j)}=\sum_{\ell=1}^n\ell^{-j}$ and $H_n=H_n^{(1)}$, with
$H_0^{(j)}=0$.

\begin{theorem}[Harmonic Bell formula for every resonant jet]\label{thm:Bell}
Fix $m\ge1$, $k>m/2$, $n=2k-m$, $r\ge1$, and $x>0$. Define
\begin{align}
 b_1&=\log x+\gamma-H_{n-1},\label{eq:b1}\\
 b_j&=(j-1)!\left((-1)^{j+1}\zeta(j)-H_{n-1}^{(j)}\right),
                  \qquad j\ge2.\label{eq:bj}
\end{align}
Then the individual late coefficient has the exact order-derivative formula
\begin{equation}\label{eq:Bell-jet}
 \left.\partial_s^r\!\left(A_k(s)x^{s-2k}\right)\right|_{s=m}
  =2(-1)^{m+1}\eta(2k)\Gamma(n)x^{-n}
       r\Bell_{r-1}(b_1,\ldots,b_{r-1}).
\end{equation}
The Stokes jump of the full $r$th order jet is obtained by differentiating
$-2\ii\sin(\pi s)T_s$ exactly $r$ times at $s=m$.
\end{theorem}
\begin{proof}
Divide \eqref{eq:gamma-detuning}, after multiplication by $x^{-n+\epsilon}$,
by $2(-1)^{m+1}\eta(2k)\Gamma(n)x^{-n}$. The result is
\[
 \epsilon\,\frac{\sin\pi\epsilon}{\pi\epsilon}
        \frac{\Gamma(n-\epsilon)}{\Gamma(n)}x^\epsilon.
\]
The logarithm of the factor multiplying $\epsilon$ is analytic near zero.
Its first derivative is $\log x-\psi(n)=b_1$.
For $j\ge2$, use
\[
 \psi^{(j-1)}(n)=(-1)^j(j-1)!
               \bigl(\zeta(j)-H_{n-1}^{(j)}\bigr)
\]
and
$\log(\sin\pi\epsilon/(\pi\epsilon))
 =-\sum_{\ell\ge1}\zeta(2\ell)\epsilon^{2\ell}/\ell$.
The $j$th derivative is precisely \eqref{eq:bj}.
The Bell generating identity now gives \eqref{eq:Bell-jet}, since
$r!/(r-1)!=r$. The full jump statement follows from the locally uniform
order differentiation already proved in Theorem~\ref{thm:lattice}.
\end{proof}

For $r=1,2,3$ the respective Bell factors are
$1$, $2b_1$, and $3(b_1^2+b_2)$. The sign
$b_2=-\zeta(2)-H_{n-1}^{(2)}$ is worth noting: it includes the contribution
of the sine factor, not only the Gamma derivative.

\subsection{A noncommuting resonance and large-order limit}

\begin{corollary}[Uniform coefficient detuning]\label{cor:detuning}
With $n=2k-m\to\infty$, and the value at $\epsilon=0$ understood by
continuity,
\begin{equation}\label{eq:detuning-ratio}
 \frac{A_k(m+\epsilon)}{\epsilon A_k'(m)}
  =\frac{\sin\pi\epsilon}{\pi\epsilon}
          \frac{\Gamma(n-\epsilon)}{\Gamma(n)}
  =\frac{\sin\pi\epsilon}{\pi\epsilon}\,
          n^{-\epsilon}\bigl(1+O(n^{-1})\bigr).
\end{equation}
The last estimate is uniform for $\epsilon$ in any fixed compact subset of
$\C$. If $\epsilon\to0$ and $\epsilon\log n\to\lambda\in\C$, then
\begin{equation}\label{eq:detuning-limit}
 \frac{A_k(m+\epsilon)}{\epsilon A_k'(m)}\longrightarrow\ee^{-\lambda}.
\end{equation}
For the coefficient including $x^{s-2k}$, the right side of
\eqref{eq:detuning-ratio} is multiplied by $x^\epsilon$.
\end{corollary}
\begin{proof}
The first identity is \eqref{eq:gamma-detuning}. Uniform Stirling expansion
of the Gamma quotient~\cite{DLMFratio}, with $n$ tending to infinity on the positive real
axis and bounded $\epsilon$, gives the second. The limit follows from
$n^{-\epsilon}=\exp(-\epsilon\log n)$ and the removable value $1$ of the
sinc factor. This is a coefficient-level uniform statement; it does not
by itself assert uniform optimal truncation of the full function.
\end{proof}

\section{Periodic residues: finite digamma and cyclotomic coordinates}\label{sec:periodic}

The lattice kernel admits an arithmetic extension that directly meets the
finite-Fourier polylogarithm--polygamma bridge of the repository
\cite{RepoBridge}. The finite Fourier transform itself is classical;
the point here is to apply it to a \emph{meromorphic Borel kernel} and its
entire exponentially small jump, including the zero Fourier coefficient.

Let $c=(c_n)_{n\ge1}$ have period $d\ge1$, with complex entries. Set
\begin{align}
 L_c(j)&=\sum_{n\ge1}\frac{c_n}{n^j}
       =d^{-j}\sum_{a=1}^dc_a\zeta(j,a/d),\qquad j\ge2,
                                                  \label{eq:Lc}\\
 T_{m,c}(x)&=\sum_{n\ge1}\frac{c_n\ee^{-nx}}{n^m},\qquad x>0.
                                                  \label{eq:Tmc}
\end{align}
No condition $\sum_ac_a=0$ is imposed.

\begin{theorem}[Periodic Borel-kernel dictionary]\label{thm:periodic}
For $m\ge1$, the normally convergent meromorphic sum
\begin{equation}\label{eq:Bmc-def}
 B_{m,c}(\xi)=\sum_{n\ge1}\frac{c_n}{n^m}
         \left(\frac1{n+\xi}-\frac{(-1)^m}{n-\xi}\right)
\end{equation}
has the finite digamma representation
\begin{align}\label{eq:Bmc-digamma}
 B_{m,c}(\xi)=\frac{(-1)^m}{\xi^m}\Bigg\{
 &\frac1d\sum_{a=1}^dc_a
    \left[\psi\!\left(\frac{a-\xi}{d}\right)
          -\psi\!\left(\frac{a+\xi}{d}\right)\right]\notag\\
 &+2\sum_{k=1}^{\lfloor m/2\rfloor}L_c(2k)\xi^{2k-1}\Bigg\}.
\end{align}
The singularity at zero is removable. The positive-pole residues and the
lateral Laplace jump are
\begin{align}
 \Res_{\xi=n}B_{m,c}(\xi)&=(-1)^m\frac{c_n}{n^m},\label{eq:periodic-res}\\
 \int_+\ee^{-x\xi}B_{m,c}(\xi)\dd\xi
 -\int_-\ee^{-x\xi}B_{m,c}(\xi)\dd\xi
 &=-2\pi\ii(-1)^mT_{m,c}(x).\label{eq:periodic-jump}
\end{align}
Writing $\omega=\ee^{2\pi\ii/d}$ and
$\widehat c_p=d^{-1}\sum_{a=1}^dc_a\omega^{-pa}$ gives the finite identity
\begin{equation}\label{eq:periodic-DFT}
 T_{m,c}(x)=\sum_{p=0}^{d-1}\widehat c_p\Li_m(\omega^p\ee^{-x}).
\end{equation}
The constants in \eqref{eq:Bmc-digamma} also have finite polygamma
coordinates:
\begin{equation}\label{eq:periodic-polygamma}
 L_c(2k)=\frac1{d^{2k}(2k-1)!}
           \sum_{a=1}^dc_a\psi^{(2k-1)}(a/d).
\end{equation}
\end{theorem}
\begin{proof}
Introduce the absolutely convergent, subtracted resolvent
\begin{equation}\label{eq:Sc}
 S_c(z)=\sum_{n\ge1}c_n\left(\frac1{n+z}-\frac1n\right)
       =\frac1d\sum_{a=1}^dc_a
          \left[\psi(a/d)-\psi((a+z)/d)\right].
\end{equation}
The equality follows by grouping the standard digamma difference series
in residue classes. Its convergence does not require a zero mean.
For
$R_m(z)=\sum_{n\ge1}c_n/[n^m(n+z)]$ one has
$R_1=-S_c/z$ and, for $m\ge2$,
$R_m=(L_c(m)-R_{m-1})/z$. Hence
\[
 R_m(z)=\sum_{j=2}^m\frac{(-1)^{m-j}L_c(j)}{z^{m-j+1}}
                   +\frac{(-1)^m S_c(z)}{z^m}.
\]
Now form $R_m(\xi)-(-1)^mR_m(-\xi)$. The odd-$j$ terms cancel and
\eqref{eq:Bmc-digamma} follows from \eqref{eq:Sc}.
The local Taylor expansion of \eqref{eq:Bmc-def} is
\begin{equation}\label{eq:Bmc-series}
 B_{m,c}(\xi)=2(-1)^{m+1}
         \sum_{k>m/2}L_c(2k)\xi^{2k-m-1},\qquad |\xi|<1,
\end{equation}
which proves removability. This is also the resonant order-jet kernel of
\eqref{eq:bs} with $(-1)^{n-1}$ replaced by $c_n$.

The residue in \eqref{eq:periodic-res} is immediate from its single pole
term. On tilted rays the finite digamma expression has at most
polynomial--logarithmic growth. Closing the upper-minus-lower contour in
the Laplace-decaying sector gives minus $2\pi\ii$ times the positive
residue sum. One can take outer radii avoiding the integer poles; the
exponential factor makes the outer arcs and residue tails tend to zero.
This proves \eqref{eq:periodic-jump}.
Finally, $c_n=\sum_p\widehat c_p\omega^{pn}$ and absolute convergence of
\eqref{eq:Tmc} give \eqref{eq:periodic-DFT}. Grouping the defining Dirichlet
series gives \eqref{eq:Lc}; the standard Hurwitz-zeta/polygamma identity
gives \eqref{eq:periodic-polygamma}.
\end{proof}

\begin{remark}[A nonzero-mean example]\label{rmk:zero-mode}
Take $c_n=1$ for $n\equiv1\pmod3$ and zero otherwise. Then
\[
 B_{1,c}(\xi)=-\frac1{3\xi}
       \left[\psi\!\left(\frac{1-\xi}{3}\right)
             -\psi\!\left(\frac{1+\xi}{3}\right)\right],
\]
and
\[
 T_{m,c}(x)=\frac13\sum_{p=0}^2
                  \omega^{-p}\Li_m(\omega^p\ee^{-x}).
\]
The $p=0$ term is $\Li_m(\ee^{-x})/3$ and cannot be dropped.
For the alternating sequence $c_n=(-1)^{n-1}$, the theorem reduces to
Theorem~\ref{thm:Gamma-kernel}. Thus one and the same residue dictionary
has Gamma, polygamma, Hurwitz-zeta, and cyclotomic-polylogarithm forms.
\end{remark}

\section{Certified inversion of the integer-order completion}\label{sec:general-inverse}

At integer order, a finite algebraic core is known exactly. Retaining it
unexpanded avoids the artificial proliferation of algebraic terms that
would result from first approximating its inverse.
Fix $m\ge1$ and $u>0$, put $y=P_m(u)$, and regard $q$ as an independent
complex parameter. The inverse displacement is the germ $\delta(q)$ with
$\delta(0)=0$ satisfying
\begin{equation}\label{eq:general-inverse-equation}
 P_m(u+\delta)-P_m(u)+(-1)^m\Li_m(-q\ee^{-\delta})=0.
\end{equation}
The physical substitution is $q=\ee^{-u}$.
Let
\begin{equation}\label{eq:Qcore}
 Q_{m,u}(z)=\frac{P_m(u+z)-P_m(u)}z,
 \qquad Q_{m,u}(0)=P_m'(u)>0,
\end{equation}
and define the finite composition sum
\begin{equation}\label{eq:composition}
 C^{(m)}_{N,k}=\coeff{t}{N}\Li_m(t)^k
   =\sum_{\substack{n_1+\cdots+n_k=N\\n_j\ge1}}
                     \frac1{(n_1\cdots n_k)^m}.
\end{equation}

\begin{theorem}[Exact sectors and an analytic certificate]\label{thm:general-inverse}
The germ $\delta(q)=\sum_{N\ge1}d_N^{(m)}(u)q^N$ has coefficients
\begin{equation}\label{eq:general-sectors}
 d_N^{(m)}(u)=(-1)^N\sum_{k=1}^N
  \frac{(-1)^{(m+1)k}C^{(m)}_{N,k}}{k!}
  \left.\frac{\mathrm d^{k-1}}{\mathrm dz^{k-1}}
       \left(\ee^{-Nz}Q_{m,u}(z)^{-k}\right)\right|_{z=0}.
\end{equation}
For a convergence certificate, choose $r>0$ such that
$P_m(u+z)-P_m(u)$ has exactly one zero in $|z|\le r$, namely the simple
zero at zero, and no boundary zero. Set
\[
 L=\min_{|z|=r}|P_m(u+z)-P_m(u)|.
\]
If $\rho\ee^r<1$ and
\begin{equation}\label{eq:certificate}
 \frac{\rho\ee^r}{1-\rho\ee^r}<L,
\end{equation}
then $\delta$ is analytic for $|q|\le\rho$, satisfies $|\delta|<r$, and,
for $|q|<\rho$,
\begin{equation}\label{eq:sector-remainder}
 \left|\delta(q)-\sum_{n=1}^N d_n^{(m)}(u)q^n\right|
   \le r\frac{(|q|/\rho)^{N+1}}{1-|q|/\rho}.
\end{equation}
For each fixed $m\ge2$, $r=1/2$ and $\rho=1/4$ satisfy the certificate
for all sufficiently large $u$. Thus the physical inverse has a genuinely
convergent exponential-sector representation at large positive core
parameter.
\end{theorem}
\begin{proof}
The implicit derivative in $\delta$ at $(q,\delta)=(0,0)$ is $P_m'(u)>0$,
so the analytic germ exists and is unique. Insert an auxiliary formal
parameter $\lambda$ into
\[
 z=\lambda\,
   \frac{(-1)^{m+1}\Li_m(-q\ee^{-z})}{Q_{m,u}(z)}.
\]
Lagrange inversion, first in $\lambda$, gives the $k$th derivative term in
\eqref{eq:general-sectors}. The coefficient of $q^N$ in the $k$th power
of the polylogarithm is $(-1)^NC^{(m)}_{N,k}\ee^{-Nz}$, and it is zero
for $k>N$. Therefore setting $\lambda=1$ is legitimate coefficientwise
in the $q$-adic ring and gives the finite formula. This is a specialization
of the general contour reversion mechanism already developed in
\cite{RepoTransport}.

For $|z|\le r$, $|q|\le\rho$, the perturbation is analytic and has
modulus at most $\rho\ee^r/(1-\rho\ee^r)$, since $m\ge1$.
Rouch\'e's theorem and \eqref{eq:certificate} give exactly one zero,
counting multiplicity, of \eqref{eq:general-inverse-equation} in the disk.
It is simple, hence depends holomorphically on $q$. Uniqueness glues the
local solutions on the parameter disk. Cauchy's estimate gives
$|d_n^{(m)}|\le r\rho^{-n}$ and summing a geometric tail proves
\eqref{eq:sector-remainder}.

For fixed $r$ and $m\ge2$,
\[
 P_m(u+z)-P_m(u)=P_m'(u)z\,(1+O(u^{-1}))
\]
uniformly on $|z|\le r$, and
$P_m'(u)\sim u^{m-1}/(m-1)!$. Comparison with this linear function
certifies the single core zero and makes $L\to\infty$.
The stated fixed $r,\rho$ then work. At $q=\ee^{-u}$ the root gives the
inverse of $F_m$ by \eqref{eq:integer-exact}; it is the real branch near
$u$. Its local monotonicity also follows from
$F_m'=F_{m-1}>0$ for $m\ge2$ on the real axis.
\end{proof}

For $m=1$ the solution can be found without a certificate:
\begin{equation}\label{eq:m1inverse}
 \delta(q)=\log(1-q)=-\sum_{N\ge1}\frac{q^N}{N},\qquad |q|<1.
\end{equation}
This checks both the signs and the normalization in
\eqref{eq:general-sectors}. The general certificate is sufficient; it is
not asserted to give an exact radius. The next sections replace it by a
sharp result for $m=2$.

\section{Dilogarithm sectors are positive Laurent polynomials}\label{sec:dilog}

For $m=2$, the exact-core equation is
\begin{equation}\label{eq:dilog-delta}
 u\delta+\frac{\delta^2}{2}+\Li_2(-q\ee^{-\delta})=0,
 \qquad \delta(0)=0.
\end{equation}
Reflecting both variables gives a positive equation. Namely, let
\begin{equation}\label{eq:vdef}
 v_u(r)=-\delta_u(-r).
\end{equation}
Then
\begin{equation}\label{eq:positive-equation}
 uv-\frac{v^2}{2}=\Li_2(r\ee^v),\qquad
 v=\Phi_u(r,v):=\frac{v^2}{2u}+\frac{\Li_2(r\ee^v)}u.
\end{equation}

\begin{theorem}[Positive finite formula]\label{thm:positive}
For every $u>0$, write
$v_u(r)=\sum_{N\ge1}a_N(u)r^N$. With $C_{N,k}=C^{(2)}_{N,k}$,
\begin{equation}\label{eq:positive-coeff}
 a_N(u)=\sum_{k=1}^N\frac{C_{N,k}}k
       \sum_{j=0}^{k-1}
       \binom{k+j-1}{j}
       \frac{N^{k-1-j}}{(k-1-j)!\,2^j\,u^{k+j}}.
\end{equation}
Every coefficient of every power $u^{-1},\ldots,u^{-(2N-1)}$ is a
strictly positive rational number, and no other powers occur. The two
extreme coefficients are
\begin{equation}\label{eq:extremes}
 [u^{-1}]a_N(u)=\frac1{N^2},\qquad
 [u^{-(2N-1)}]a_N(u)=\frac{\operatorname{Cat}_{N-1}}{2^{N-1}},
 \quad \operatorname{Cat}_{j}=\frac1{j+1}\binom{2j}{j}.
\end{equation}
Consequently $d_N^{(2)}(u)=(-1)^{N-1}a_N(u)$.
\end{theorem}
\begin{proof}
Here $Q_{2,u}(z)=u+z/2$. In the $(k-1)$st derivative of
$\ee^{-Nz}(u+z/2)^{-k}$ every Taylor coefficient contributing to that
order has sign $(-1)^{k-1}$. Expanding the exponential and the negative
binomial power in \eqref{eq:general-sectors} therefore gives
\eqref{eq:positive-coeff}, with total sign $(-1)^{N-1}$.

The rational numbers $C_{N,k}$ are positive whenever $1\le k\le N$.
For any integer $1\le p\le2N-1$, one can choose $k\le N$ and
$0\le j<k$ with $p=k+j$, so every indicated power occurs positively.
The power $u^{-1}$ comes only from $k=1,j=0$, where $C_{N,1}=N^{-2}$.
The highest power comes only from $k=N,j=N-1$, where $C_{N,N}=1$.
Its coefficient is
$\binom{2N-2}{N-1}/(N2^{N-1})$, which is the Catalan expression.
\end{proof}

The first three coefficients are
\begin{align}
 a_1(u)&=\frac1u,\notag\\
 a_2(u)&=\frac1{4u}+\frac1{u^2}+\frac1{2u^3},\label{eq:first-coeff}\\
 a_3(u)&=\frac1{9u}+\frac3{4u^2}+\frac7{4u^3}
                   +\frac3{2u^4}+\frac1{2u^5}.\notag
\end{align}
A separate exact recurrence, obtained by iterating the positive map
$\Phi_u$, verifies the finite formula independently of Lagrange inversion
in the accompanying program. The positivity, rather than merely the
existence of a coefficient formula, is the key to the global radius proof.

\section{The exact radius and its unique dominant singularity}\label{sec:radius}

\begin{theorem}[Sharp dilogarithm reversion radius]\label{thm:radius}
Fix $u>0$. There is a unique $t=t_u\in(0,1)$ satisfying
\begin{equation}\label{eq:characteristic-t}
 u^2=2\Li_2(t)+\log^2(1-t).
\end{equation}
Set
\begin{equation}\label{eq:radius-parameters}
 w=-\log(1-t),\qquad v_c=u-w,\qquad
 R=R(u)=t\ee^{w-u}=\ee^{-u}\frac{t}{1-t}.
\end{equation}
Then $0<w<u$, $v_c>0$, and the Taylor series of $v_u$ has exact radius
$R$. The point $r=R$ is its unique singularity on the circle $|r|=R$.
In a slit neighborhood of $R$,
\begin{equation}\label{eq:sqrt-expansion}
 v_u(r)=v_c-\sqrt{2w(1-t)}\sqrt{1-r/R}+O(1-r/R),
\end{equation}
with a convergent Puiseux expansion in $\sqrt{1-r/R}$.
The coefficients satisfy
\begin{equation}\label{eq:large-sector}
 a_N(u)=\sqrt{\frac{w(1-t)}{2\pi}}\,
           R^{-N}N^{-3/2}\bigl(1+O(N^{-1})\bigr).
\end{equation}
All constants and error terms here are for fixed $u$. Equivalently, the
unique dominant singularity of $\delta_u(q)$ is at $q=-R$.
\end{theorem}

\begin{proof}
\emph{The real characteristic point.}
The function $2\Li_2(t)+\Li_1(t)^2$ is strictly increasing from zero to
infinity on $(0,1)$, since its derivative is positive and
$\Li_1(t)=-\log(1-t)\to\infty$. Thus \eqref{eq:characteristic-t} has a
unique solution. It implies $w^2<u^2$ and $v_c>0$.
For $0\le t\le t_u$ define
\begin{equation}\label{eq:parametric}
 v(t)=u-\sqrt{u^2-2\Li_2(t)},\qquad r(t)=t\ee^{-v(t)}.
\end{equation}
The radicand is positive on this whole interval. Differentiation gives
\begin{equation}\label{eq:param-derivative}
 \frac{\mathrm d}{\mathrm dt}\log r(t)
   =\frac1t\left(1-
         \frac{\Li_1(t)}{\sqrt{u^2-2\Li_2(t)}}\right).
\end{equation}
It is positive for $0<t<t_u$ and vanishes at $t_u$.
Therefore \eqref{eq:parametric} supplies a real analytic solution for
$0\le r<R$, tending to $v_c$ at $R$.

\emph{No smaller convergence radius.}
Starting from $v_0=0$, iterate $v_{j+1}=\Phi_u(r,v_j)$ as a formal power
series. Coefficients are nonnegative and increase coefficientwise. The
coefficient of each fixed degree stabilizes, because
$\partial_v\Phi_u(0,0)=0$ makes the difference between two successive
iterates gain at least one order in $r$. The limit is the unique formal
solution. For fixed real $0\le r<R$, the iterates are bounded above by
the real solution in \eqref{eq:parametric}, since $\Phi_u$ preserves
order on the relevant positive interval. Monotone convergence of the
nonnegative coefficients shows that the full series converges for every
such $r$, hence absolutely for $|r|<R$. Its sum is the real branch just
constructed, by uniqueness near zero and analytic continuation.
Moreover, monotone convergence as $r\uparrow R$ gives
\begin{equation}\label{eq:boundary-sum}
 \sum_{N\ge1}a_N(u)R^N=v_c<\infty.
\end{equation}

\emph{The singularity is a genuine square-root fold.}
At $(R,v_c)$ the defining equation $v=\Phi_u(r,v)$ satisfies
\begin{equation}\label{eq:fold-derivatives}
 \Phi_v=1,\qquad R\Phi_r=\frac wu,\qquad
 \Phi_{vv}=\frac{1+\Li_0(t)}u=\frac1{u(1-t)}.
\end{equation}
Both nonzero derivatives are positive. The analytic implicit function
theorem applied after solving for $r$ as a function of $v$, followed by
factoring its double zero at $v_c$, gives a convergent square-root
inverse. In more detail, write $r=R-\alpha(v-v_c)^2+
O((v-v_c)^3)$ with $\alpha>0$; the analytic square root of the nonvanishing
factor reduces this to an ordinary analytic inverse in
$\sqrt{R-r}$. Balancing the first nonzero terms in $v-\Phi_u(r,v)=0$
gives the coefficient $\sqrt{2R\Phi_r/\Phi_{vv}}=
\sqrt{2w(1-t)}$. The negative sign follows from $v(r)<v_c$ for $r<R$.
This proves \eqref{eq:sqrt-expansion} and shows that the radius is not
larger than $R$.

\emph{Uniqueness on the dominant circle.}
By \eqref{eq:boundary-sum}, the series converges uniformly and absolutely
on the closed disk. For $|r|=R$, $r\ne R$, strict triangle inequality gives
$|v(r)|<v_c$, since every coefficient $a_N$ is positive and the positive
support includes all $N\ge1$. Consequently
\[
 |r\ee^{v(r)}|=R\ee^{\Re v(r)}<R\ee^{v_c}=t_u<1.
\]
The defining polylogarithm is analytic there, and
\[
 |\Phi_v(r,v(r))|
 \le\frac{|v(r)|+\Li_1(R\ee^{\Re v(r)})}u
 <\frac{v_c+w}u=1.
\]
The implicit function theorem thus extends $v$ analytically through every
other point of the circle. A finite covering of its compact complement
and the local square-root chart at $R$ give a standard dented-disk
($\Delta$-domain) continuation about the unique dominant singularity.

\emph{Coefficient asymptotics.}
The elementary singular term satisfies
\[
 [r^N](1-r/R)^{1/2}
   =-\frac{R^{-N}}{2\sqrt\pi}\,N^{-3/2}(1+O(N^{-1})).
\]
The convergent Puiseux expansion and the preceding $\Delta$-continuation
permit termwise singularity transfer~\cite{FO}. Analytic integer-power
terms contribute no competing singular term; the next half-integer
power has coefficient $O(R^{-N}N^{-5/2})$. This proves
\eqref{eq:large-sector}. The relation \eqref{eq:vdef} transports the
singularity to $q=-R$.
\end{proof}

\begin{remark}[Why a real fold alone would not suffice]
Solving the two real characteristic equations would provide a candidate
radius, but would not exclude a closer complex singularity. The positive
iteration and the strict boundary derivative estimate in the proof are
what certify the \emph{exact} radius and its unique dominant point.
\end{remark}

\section{The sharp physical threshold and an inverse identity}\label{sec:physical}

\begin{theorem}[Necessary and sufficient physical convergence]\label{thm:threshold}
Let $y>\pi^2/6$ and $u$ be given by \eqref{eq:u-y}. Then
\begin{equation}\label{eq:physical-series}
 F_2^{-1}(y)=u+\sum_{N\ge1}(-1)^{N-1}a_N(u)\ee^{-Nu}
\end{equation}
is an absolutely convergent identity if and only if
\begin{equation}\label{eq:physical-threshold}
 u\ge\frac\pi{\sqrt6},\qquad\text{equivalently}\qquad
 y\ge\frac{\pi^2}{4}.
\end{equation}
At equality convergence is absolute despite evaluation on the Taylor
circle. Below the threshold, the terms on the right do not tend to zero.
\end{theorem}
\begin{proof}
The physical parameter is $q=\ee^{-u}$. By \eqref{eq:radius-parameters},
$q\le R$ is equivalent to $t/(1-t)\ge1$, or $t\ge1/2$.
The characteristic function in \eqref{eq:characteristic-t} is strictly
increasing. The classical reflection identity gives
\[
 \Li_2(1/2)=\frac{\pi^2}{12}-\frac{\log^2 2}{2},\qquad
 2\Li_2(1/2)+\log^2 2=\frac{\pi^2}{6}.
\]
Thus $q\le R$ is exactly \eqref{eq:physical-threshold}.
For $q<R$, the series is the analytic germ from \eqref{eq:dilog-delta}.
For every real $q>0$, that equation has a unique positive solution:
its left side is negative at $\delta=0$, tends to infinity, and has
strictly positive $\delta$ derivative
\[
 u+\delta+\log(1+q\ee^{-\delta})>0.
\]
Consequently the analytic germ is this real root wherever its Taylor
series converges, and \eqref{eq:integer-exact} identifies the root with
the required inverse at $q=\ee^{-u}$.

At $q=R$, \eqref{eq:boundary-sum} or \eqref{eq:large-sector} gives absolute
convergence. Continuity extends the same defining equation and hence the
same inverse identity to the endpoint. If $q>R$, \eqref{eq:large-sector}
gives
$a_Nq^N\sim C(u)(q/R)^NN^{-3/2}$, with $C(u)>0$; terms do not tend to
zero. Finally $y=u^2/2+\pi^2/6$ converts the threshold in $u$ to that in
$y$.
\end{proof}

The derivative $F_2'(x)=\log(1+\ee^x)>0$ and the limits at $\pm\infty$
show that $F_2$ maps $\R$ bijectively onto $(0,\infty)$. Therefore the
failure below \eqref{eq:physical-threshold} is solely a failure of this
expansion, not of the real inverse.

\begin{corollary}[An exact self-referential radius identity]\label{cor:radius-inverse}
For every $u>0$,
\begin{equation}\label{eq:radius-inverse}
 R(u)=\exp\bigl(F_2^{-1}(u^2/2)-u\bigr).
\end{equation}
Also $0<R(u)<1$, $R(u)\sim u^2/2$ as $u\downarrow0$, and
\begin{equation}\label{eq:radius-large-u}
 \log R(u)=-\frac{\pi^2}{6u}+O(u^{-3})\qquad(u\to\infty).
\end{equation}
\end{corollary}
\begin{proof}
Let $a=t/(1-t)=\ee^w-1$. The classical Landen identity, obtained also by
differentiating both sides and matching at $t=0$, is
\[
 -\Li_2\!\left(-\frac{t}{1-t}\right)
     =\Li_2(t)+\frac12\log^2(1-t).
\]
The characteristic equation therefore gives $F_2(\log a)=u^2/2$.
Since $R=\ee^{-u}a$, \eqref{eq:radius-inverse} follows.
It is an exact functional relation, not an independent shortcut for
computing an inverse one has not yet evaluated.

The inequalities $t<1$ and $w<u$ give $R=t\ee^{w-u}<1$.
In $w$ coordinates the characteristic equation is
\begin{equation}\label{eq:characteristic-w}
 u^2=w^2+2\Li_2(1-\ee^{-w}).
\end{equation}
Its right side has derivative $2w/(1-\ee^{-w})$, with limiting value $2$
at zero. Hence $w\sim u^2/2$ and $R=\ee^{-u}(\ee^w-1)\sim u^2/2$.
For large $u$, $w\to\infty$, and the reflection identity gives
$\Li_2(1-\ee^{-w})=\pi^2/6+O(w\ee^{-w})$.
Thus
$w=u-\pi^2/(6u)+O(u^{-3})+O(\ee^{-w})$.
Using $\log R=w-u+\log(1-\ee^{-w})$ proves
\eqref{eq:radius-large-u}.
\end{proof}

\section{A quadratic polylogarithm-feedback extension}\label{sec:feedback}

The sharp-radius argument depends more on positivity than on the special
functional identities of the dilogarithm. Those identities are what turn
the general radius criterion into the especially simple threshold
$\pi^2/4$.

\begin{theorem}[Interior folds for quadratic polylogarithm feedback]\label{thm:feedback}
Let $s>0$, $u>0$, and let $v$ be the germ with $v(0)=0$ satisfying
\begin{equation}\label{eq:feedback}
 uv-\frac{v^2}{2}=\Li_s(r\ee^v).
\end{equation}
Suppose there is $t\in(0,1)$ with
\begin{equation}\label{eq:s-characteristic}
 u^2=2\Li_s(t)+\Li_{s-1}(t)^2.
\end{equation}
This condition holds for every $u>0$ when $0<s\le2$. For $s>2$ it holds
exactly when
\begin{equation}\label{eq:Uc}
 u<U_c(s):=\sqrt{2\zeta(s)+\zeta(s-1)^2}.
\end{equation}
The solution $t$ is unique. Put
\begin{equation}\label{eq:s-radius}
 w=\Li_{s-1}(t),\quad v_c=u-w,\quad R=t\ee^{w-u}.
\end{equation}
Then the Taylor series of $v$ has positive coefficients, exact radius $R$,
and unique dominant singularity $R$. Its coefficient asymptotic is
\begin{equation}\label{eq:s-asymptotic}
 [r^N]v(r)\sim
 \sqrt{\frac{w}{2\pi\bigl(1+\Li_{s-2}(t)\bigr)}}\,
       R^{-N}N^{-3/2}.
\end{equation}
The positive finite formula \eqref{eq:positive-coeff} remains valid after
replacing $C_{N,k}$ by $\coeff{z}{N}\Li_s(z)^k$. Its coefficients in
$u^{-1}$ are then positive real numbers; the extreme coefficients are
$N^{-s}$ and $\operatorname{Cat}_{N-1}/2^{N-1}$.
\end{theorem}
\begin{proof}
For real $s$, each coefficient $n^{-s}$ is positive. All derivatives with
respect to $t$ of a polylogarithm in the open interval $(0,1)$ are
therefore positive. In particular,
$2\Li_s(t)+\Li_{s-1}(t)^2$ is strictly increasing. Its limit at $t=1$
is infinite if $s\le2$, and is $U_c(s)^2$ if $s>2$. This proves the
existence and uniqueness assertions.

Use the same parameterization as \eqref{eq:parametric}, with $\Li_s$
in place of $\Li_2$. Its logarithmic $r$ derivative is
$t^{-1}(1-\Li_{s-1}(t)/\sqrt{u^2-2\Li_s(t)})$, positive before the
characteristic point. Iterating
$\Phi(r,v)=v^2/(2u)+\Li_s(r\ee^v)/u$ gives nonnegative formal
coefficients and bounds their sum up to $R$ exactly as in the proof of
Theorem~\ref{thm:radius}. The Lagrange calculation in
Theorem~\ref{thm:positive} uses only the quadratic core, so its positive
finite formula applies unchanged.

At the fold,
\[
 R\Phi_r=w/u,\qquad \Phi_{vv}=(1+\Li_{s-2}(t))/u>0.
\]
Hence the square-root coefficient is
$\sqrt{2w/(1+\Li_{s-2}(t))}$.
Every coefficient is positive. On another point of the circle,
$|v|<v_c$, $|r\ee^v|<t$, and
$|\Phi_v|\le(|v|+\Li_{s-1}(|r\ee^v|))/u<1$.
This proves unique dominance and the same $\Delta$ continuation as before.
The square-root transfer gives \eqref{eq:s-asymptotic}.
\end{proof}

For $0<s\le2$, the boundary between convergence and divergence on the
curve $r=-\ee^{-u}$ can likewise be stated exactly. The unique
$t_*\in(0,1)$ solving
\begin{equation}\label{eq:general-physical}
 \log t_*+\Li_{s-1}(t_*)=0
\end{equation}
gives $u_*^2=2\Li_s(t_*)+\Li_{s-1}(t_*)^2$.
Convergence holds precisely for $u\ge u_*$, including absolute convergence
at the boundary. This follows from $\ee^{-u}/R=\ee^{-w}/t$ and monotonicity
of both sides of the characteristic parametrization. For $s=2$,
$t_*=1/2$. For other values of $s$, \eqref{eq:feedback} is a quadratic
feedback model, \emph{not} the exact-core inversion equation of $F_s$;
that identification is special to the dilogarithm case.

\section{An audit safeguard: arithmetic averaging and inversion}\label{sec:averaging}

The repository source~\cite{RepoTransport} already explicitly studies the
nonlinear real-averaging defect. The calculation below is an illustrative
restatement and integration safeguard, not a newly discovered phenomenon.
The exact identity \eqref{eq:average} uses an arithmetic average of
lateral realizations. It must not be propagated through nonlinear
inversion as though averaging were a homomorphism. The following elementary
calculation isolates the obstruction.

\begin{proposition}[First nonlinear averaging defect]\label{prop:average-defect}
Let $F$ be analytic and locally invertible with inverse $g$, let $T$ be
analytic near the relevant inverse point, and let $G_\pm(y;\lambda)$ be the
coherently chosen inverses of $F\pm\ii\lambda T$ for sufficiently small
$\lambda$. Then
\begin{equation}\label{eq:averaging-defect}
 \frac{G_++G_-}{2}
   =g-\frac{\lambda^2}{2}\frac{\mathrm d}{\mathrm dy}
               \bigl(g'(y)T(g(y))^2\bigr)+O(\lambda^4),
\end{equation}
locally uniformly. The inverse of the arithmetic average of the forward
functions is exactly $g$, so these operations need not commute.
\end{proposition}
\begin{proof}
The analytic perturbation inverse has the convergent local expansion
\[
 G_\pm=g+\sum_{n\ge1}\frac{(\mp\ii\lambda)^n}{n!}
             \frac{\mathrm d^{n-1}}{\mathrm dy^{n-1}}
               \bigl(g'(y)T(g(y))^n\bigr).
\]
It follows either by a local contour logarithm, as in
\cite{RepoTransport}, or by Lagrange inversion in the forward coordinate.
Averaging removes the odd terms; the $n=2$ coefficient is the displayed
one. Holomorphy in $\lambda$ gives the uniform fourth-order remainder.
\end{proof}

For $F(x)=x$ and $T(x)=\ee^{-x}$, the defect is
$+\lambda^2\ee^{-2y}+O(\lambda^4)$. This gives an explicit counterexample,
not just a warning based on formal terminology. It is relevant when the
forward polylogarithm happens to equal an arithmetic lateral average,
as at a half-integer order in \eqref{eq:average}.

This proposition is \emph{not} a counterexample to appropriately defined
\'Ecalle median summation or to a summation morphism whose hypotheses
hold. Arithmetic averaging and such nonlinear-compatible summation
procedures are different constructions. Nor was an erroneous averaging
claim found in the inspected accepted repository theorem. The proposition
is a preventive integration note, not an allegation about an unseen text.

\section{Verification, limitations, and proposed integration}\label{sec:verification}

\subsection{Exact replay and numerical diagnostics}

The distributed program \path{code/verify.py} uses Python rational
arithmetic for the inverse coefficients and SymPy for finite symbolic
checks. At the recorded run it passed \textbf{118 exact assertions}:
positivity and complete Laurent support through $N=10$, both extreme
coefficients at each of those orders, thirty coefficient comparisons
against an independent fixed-point recurrence at $u=1,2,3$, and finite
Taylor/residue checks of six Gamma kernels. These assertions are finite
regressions, not a formal proof of a universally quantified theorem.

The arbitrary-precision run used \texttt{mpmath} at 65 decimal digits.
Three complex-ray reconstructions of the first order derivative were
compared with differentiation under the defining Fermi integral, twelve
higher-jet coefficient comparisons were performed, and six periodic
kernels were checked in both digamma and convergent Hurwitz-series
coordinates. Their Fourier tails were separately checked against the
exponential Dirichlet series. The recorded absolute errors for these
comparisons were below $10^{-53}$; this is floating-point agreement,
not an interval certificate.

The radius equation was solved monotonically in its $w$ coordinate.
Table~\ref{tab:radius} displays selected results, and the full-precision
records are in \path{data/radii.csv}. In particular, a ratio larger than one
puts the physical point outside the convergence disk.
\begin{table}[htbp]\centering
\caption{Exact-radius diagnostics, rounded from 65-digit calculations.}
\label{tab:radius}
\begin{tabular}{rrrr}\toprule
$u$ & $w$ & $R(u)$ & $\ee^{-u}/R(u)$\\\midrule
$0.7$ & $0.231283506580$ & $0.129219674046$ & $3.842954313697$\\
$1$ & $0.447460986781$ & $0.207607344120$ & $1.771996278510$\\
$\pi/\sqrt6$ & $\log2$ & $0.277329255639$ & $1$\\
$2$ & $1.418886391532$ & $0.423939922809$ & $0.319232221254$\\
$4$ & $3.600624855684$ & $0.652423391327$ & $0.028073240678$\\
$8$ & $7.792136744823$ & $0.811982646785$ & $0.000413140144$\\\bottomrule
\end{tabular}
\end{table}

Table~\ref{tab:coeff-limit} checks the fixed-$u$ asymptotic constant by
positive convolution and the finite formula; these particular larger-$N$
calculations use binary64, not 65-digit arithmetic. Their slower approach
at $u=4$ is consistent with, but does not prove, the lack of uniformity
as $u\to\infty$.
\begin{table}[htbp]\centering
\caption{The normalized coefficient
$a_N(u)R(u)^NN^{3/2}/\sqrt{w(1-t)/(2\pi)}$, which tends to one.}
\label{tab:coeff-limit}
\begin{tabular}{rrr}\toprule
$N$ & $u=2$ & $u=4$\\\midrule
16 & 0.994010739 & 1.119387718\\
32 & 0.997036572 & 1.078796636\\
64 & 0.998528610 & 1.049118033\\
128 & 0.999267319 & 1.029009229\\
256 & 0.999634478 & 1.016318713\\\bottomrule
\end{tabular}
\end{table}

Direct inverse comparisons at $u=2$ and $u=4$ used up to 24 exact
coefficient polynomials. At $N=24$, the absolute inverse errors were about
$5.74\cdot10^{-16}$ and $1.72\cdot10^{-42}$ respectively. The associated
forward residuals were also checked. These are evaluations of a proved
convergent expansion; they are not the source of the convergence theorem.

\subsection{A numerical differentiation trap detected during the audit}

A preliminary check used a tiny-step numerical order derivative of
$-\Li_s(-\ee^{1.5})$ at $s=1$. With the installed \texttt{mpmath 1.3.0},
that route did not retain the requested accuracy. At 50-digit working
precision, one observed value was
\[
 0.9040505816228687763214111328125,
\]
whereas both the independent defining-integral derivative and the Borel
reconstruction gave
\[
 0.904050579573149170271101746468683814622500015191\ldots.
\]
The discrepancy is approximately $2.05\cdot10^{-9}$; merely quoting the
working precision would conceal it. This is an observed behavior of that
specific evaluation route and version, not a general theorem about the
library or a diagnosis of its internal cause. The final verification
program does \emph{not} use that route as its reference. It evaluates
\begin{equation}\label{eq:independent-jet-integral}
 J_m(x)=\frac1{\Gamma(m)}\int_0^\infty
    \frac{t^{m-1}(\log t-\psi(m))}{1+\ee^{t-x}}\dd t,
\end{equation}
which follows directly from \eqref{eq:FDintegral}. The package includes a
small reproducer and a version-specific note so the failed preliminary
check is not silently lost.

\subsection{Source audit and integration boundary}

The audit inspected selected portions of the project READMEs, the initial portion of the
polylogarithm--polygamma bridge article, the nonlinear Stokes-transport
title/abstract, proof, and countable-family/flow discussion, and the critical-Hahn
package README. It did not inspect every chapter of the large canonical
volumes or every pending incoming article. The last observed main commit
was \texttt{cbfbc3163e0d038a381808c6db8d39ea12dc0db7}; individual file
blob identifiers and reading ranges are recorded in \path{SOURCES.md}.
Some initial reads preceded that commit observation and were made on
\texttt{main}; the audit is therefore not falsely presented as an atomic
full-tree checkout.

No new mathematical error was established in the accepted repository
statements actually inspected. The proposed corrections are narrowly
scoped safeguards for this contribution: retain exponential completion
at integer resonance; retain order derivatives of vanishing coefficients;
fix the lateral sign convention; preserve the zero Fourier mode; do not
replace a sufficient certificate by an exact-radius claim; and do not
commute arithmetic averaging with inversion. Existing source warnings
about evidence versus proof are respected, not rebranded as discoveries.

Recommended placement is a new self-contained research package under
\path{Analysis/Transseries/docs/series-and-transseries/}, with a backlink
from the polylogarithm reports. The generic inverse engine should remain
attributed to its existing source. The sharper, model-specific radius and
periodic-kernel statements can be considered separately for promotion to
the canonical volumes after independent proof review. The package changes
no repository files remotely.

\section{Further research questions and concrete next steps}\label{sec:future}

\paragraph{1. The boundary regime of quadratic polylogarithm feedback.}
Theorem~\ref{thm:feedback} stops at an interior fold. For $s>2$ and
$u\ge U_c(s)$, the polylogarithm boundary at $t=1$ competes with or replaces
the fold. A useful precise target, when $u>U_c(s)$, is the candidate law
\[
 [r^N]v(r)\sim
 \frac{w_0^{s-1}}{(w_0-\zeta(s-1))^s}
       R_b^{-N}N^{-s},\qquad
 w_0=\sqrt{u^2-2\zeta(s)},\quad R_b=\ee^{-(u-w_0)}.
\]
This expression follows from formal transport of the leading
polylogarithm singular term, but is \emph{not proved here}. A proof must
establish the relevant boundary continuation, handle integer-order
logarithms, and exclude competing singularities. At $u=U_c(s)$, one expects
different exponents for $2<s<3$, a logarithmic transition at $s=3$, and a
finite-curvature regime above it. The related linear-core critical-Hahn
work in \cite{RepoCritical} should be used as a comparison source rather
than ignored.

\paragraph{2. Joint large-sector and large-core asymptotics.}
The constant in \eqref{eq:large-sector} tends rapidly to zero as
$u\to\infty$, while $R(u)\to1$. Determine a uniform law for $N,u\to\infty$,
including the transition from the fixed-$N$ behavior
$a_N(u)\sim1/(N^2u)$ to the fixed-$u$ square-root law. The parameter
$1-t_u=\ee^{-w_u}$ measures the approaching polylogarithm boundary and
suggests that the product of this quantity with the sector index deserves
particular attention; the correct full scaling should be derived, not
assumed from the separate limits.

\paragraph{3. Exact radii for higher integer-order cores.}
For $m\ge3$, \eqref{eq:general-sectors} is finite and explicit, but
$Q_{m,u}$ is no longer linear. Classify when an analogous sign pattern
survives, locate the nearest complex critical value, and determine whether
any further physical thresholds reduce to known special polylogarithm
values. The sufficient certificate in Theorem~\ref{thm:general-inverse}
provides a safe starting region but does not answer those questions.

\paragraph{4. A meromorphic--logarithmic dictionary for higher jets.}
Theorem~\ref{thm:Bell} gives every coefficient, whereas
Theorem~\ref{thm:Gamma-kernel} gives a particularly simple global Borel
kernel only for the first jet. Derive comparable global representations
for $r\ge2$, with the logarithmic singular germs at every integer made
explicit. The differentiated lattice formula is available, but an
arithmetic normal form should keep the harmonic sums and branch constants
separate and reveal which finite Gamma/polygamma coordinates suffice.

\paragraph{5. Resonance detuning coupled to nonlinear inversion.}
Equation~\eqref{eq:detuning-limit} resolves the moving order at the
coefficient level. Construct a uniform inverse chart when the order
approaches an integer while the core grows, and determine the scale at
which the restored divergent tail competes with the pre-existing
exponential inverse sectors. Coefficientwise convergence alone does not
justify commuting these limits with Borel integration and reversion.

\paragraph{6. Arithmetic cancellation of the first Borel actions.}
Periodic weights may cancel the first several residues exactly. Determine
how such cancellations change the first visible inverse sectors, the
nearest surviving Borel singularities, and optimal truncation. The finite
DFT in Theorem~\ref{thm:periodic} makes this a concrete algebraic design
problem. A nonzero mean must remain in the calculation; cancellation of
some residue classes does not imply disappearance of the zero mode.

\paragraph{7. Nonlinear-compatible real summation.}
Starting with the explicit jump \eqref{eq:jump}, construct and compare
real analytic completions whose inverse transport is compatible with a
chosen summation framework. Proposition~\ref{prop:average-defect} supplies
a low-order obstruction test. The goal is an actual compatible construction,
not a relabeling of arithmetic averaging as a median operation.

\paragraph{8. Formalization and certified computation.}
A practical first formalization target is the positive coefficient formula
and its two extreme coefficients over $\Q[u^{-1}]$, followed by the
finite-Fourier residue dictionary. The analytic radius theorem additionally
needs a positive implicit-series argument, a complex continuation API,
and a square-root transfer theorem. Separately, replace the current
floating-point Borel diagnostics by interval enclosures on tilted rays
with rigorous tail bounds. These tasks are proposed work; the present
package contains neither a Lean proof nor an interval certificate.

\appendix
\section{Reproduction guide and result map}\label{app:reproduce}

All manuscript formulas are contained in this one \texttt{article.tex};
no external image or font file is required. The retained tables are rounded
from the accompanying CSV files. The following commands reproduce the
verification and the PDF in a standard Python/TeX environment:
\begin{verbatim}
python -m pip install -r requirements.txt
python code/verify.py --output-dir replay --max-n 256
pdflatex -interaction=nonstopmode -halt-on-error article.tex
pdflatex -interaction=nonstopmode -halt-on-error article.tex
pdflatex -interaction=nonstopmode -halt-on-error article.tex
\end{verbatim}
Use a new output directory for a replay rather than overwriting historical
evidence. The optional \texttt{--skip-borel} switch is explicitly recorded
as a skipped check; it was \emph{not} used in the final retained run.

\small
\begin{longtable}{@{}>{\raggedright\arraybackslash}p{.28\textwidth}>{\raggedright\arraybackslash}p{.22\textwidth}>{\raggedright\arraybackslash}p{.43\textwidth}@{}}
\toprule
Result & Status in this article & Dependency and scope\\\midrule
\endhead
Theorem~\ref{thm:lattice} & Classical completion, explicit realization &
Defining integral, binomial series, normal convergence; no claim that the
complete Sommerfeld formula is new.\\[5pt]
Theorems~\ref{thm:Gamma-kernel} and~\ref{thm:Bell} & Derived jet identities &
Gamma reflection, analytic order differentiation, explicit endpoint terms;
conventional proofs.\\[5pt]
Theorem~\ref{thm:periodic} & Derived arithmetic kernel dictionary &
Subtracted digamma resolvent and finite DFT; zero mean not assumed.\\[5pt]
Theorem~\ref{thm:general-inverse} & Model specialization of existing tools &
Lagrange inversion and Rouch\'e certification; sufficient radius only.\\[5pt]
Theorems~\ref{thm:positive}--\ref{thm:threshold} & Main proposed contribution &
All-sector positivity, exact radius, unique dominance, sharp endpoint;
not established merely by the numerical checks.\\[5pt]
Theorem~\ref{thm:feedback} & Derived interior-fold extension &
All $0<s\le2$, or $s>2$ below the stated boundary; no boundary theorem
claimed at or above $U_c(s)$.\\[5pt]
Proposition~\ref{prop:average-defect} & Exact audit safeguard &
Elementary perturbative inverse calculation; not a refutation of median
summation theory.\\\bottomrule
\end{longtable}
\normalsize

\begingroup\small
\begin{thebibliography}{20}
\setlength{\itemsep}{2pt}
\bibitem{DLMFpoly}
NIST Digital Library of Mathematical Functions,
\emph{\S25.12, Polylogarithms}, especially the integral representation,
Jonqui\`ere expansion, dilogarithm reflection and Landen identities.
\url{https://dlmf.nist.gov/25.12}. Accessed 11 October 2026.

\bibitem{DLMFgamma}
NIST Digital Library of Mathematical Functions,
\emph{\S5.5, Functional Relations}, Gamma reflection.
\url{https://dlmf.nist.gov/5.5}. Accessed 11 October 2026.

\bibitem{DLMFratio}
NIST Digital Library of Mathematical Functions,
\emph{\S5.11, Asymptotic Expansions}, especially the digamma expansion
and Gamma-ratio formulas.
\url{https://dlmf.nist.gov/5.11}. Accessed 11 October 2026.

\bibitem{GFG}
T. M. Garoni, N. E. Frankel, and M. L. Glasser,
Complete asymptotic expansions of the Fermi--Dirac integrals,
\emph{Journal of Mathematical Physics} \textbf{42} (2001), 1860--1868.
\url{https://doi.org/10.1063/1.1350634}.
The present source audit inspected the institutional publication record
and abstract, not the complete article.

\bibitem{GST}
A. Gil, J. Segura, and N. M. Temme,
\emph{Complete asymptotic expansions for the relativistic Fermi--Dirac
integral}, arXiv:2108.11210v2, Section~3 and Lemma~3.1.
\url{https://arxiv.org/html/2108.11210v2}.

\bibitem{FO}
P. Flajolet and A. Odlyzko,
Singularity analysis of generating functions,
\emph{SIAM Journal on Discrete Mathematics} \textbf{3} (1990), 216--240.
\url{https://doi.org/10.1137/0403019}.

\bibitem{RepoPoly}
ProveIt Contributors,
\emph{Polylogarithms: project README and claim-status overview},
\path{Analysis/Polylogarithms/README.md}, in the ProveIt repository.
\url{https://github.com/VladimirReshetnikov/ProveIt}.
Reading scope and observed revision are recorded in \path{SOURCES.md}.

\bibitem{RepoBridge}
V. Reshetnikov,
\emph{The Polylogarithm--Polygamma Bridge: An Exact Finite Fourier
Transform, Clausen-Type Laws at Rational Arguments, and the Layers Below
the Digamma}, 11 June 2026, historical source retained at
\path{Analysis/Polylogarithms/docs/articles/polylog-polygamma-bridge.tex}.
The current project ledger, not the historical abstract alone, governs
accepted claim status.

\bibitem{RepoTrans}
ProveIt Contributors,
\emph{Series and transseries: canonical volume and research-package index},
\path{Analysis/Transseries/docs/series-and-transseries/README.md},
ProveIt repository, inspected 11 October 2026.

\bibitem{RepoTransport}
ProveIt research package,
\emph{Nonlinear Stokes Transport under Logarithmic Inversion},
\path{Analysis/Transseries/docs/series-and-transseries/Nonlinear_Stokes_Transport_Logarithmic_Inversion/article.tex}.
Inspected source lines 1--110 and 180--370; blob
\texttt{1e8b24325554d1f1c94c29f21b794ba62feff755}.
This is the source for the general contour inverse-transport architecture,
its countable-family qualification, and the distinction between a
certificate and a formal formula.

\bibitem{RepoCritical}
ProveIt research package,
\emph{Beyond Finite-Action Folds: Critical Hahn Transseries, Stable Sector
Asymptotics, and Sharp Action Budgets}, README at
\path{Analysis/Transseries/docs/series-and-transseries/Critical_Hahn_Transseries_Beyond_Finite_Action_Folds/README.md}.
Inspected at commit
\texttt{cbfbc3163e0d038a381808c6db8d39ea12dc0db7}; README blob
\texttt{7d138677987dd1cc6dbfdf5df1aa564bfd08d5dd}.
Its linear-core model and editorial cross-references are comparison
material, not analytic results audited in full here.
\end{thebibliography}
\endgroup
\end{document}
